\chapter*{Abstract}\label{chap:abstract}
\phantomsection
\addcontentsline{toc}{chapter}{Abstract}
\medskip
\noindent

This project aims to develop a prototypical slicer for non-planar, robot-assisted additive manufacturing.
Conventional slicers create purely horizontal layers when slicing objects.
This method brings inherent disadvantages, such as direction-dependent component properties as well as the need for support structures when trying to print overhangs.

To overcome these limitations, this research project aims to utilize the six degrees of freedom of an industrial robot.
Using these, the tool orientation and path planning can be adapted to better fit the object's geometry.

To solve the software-related challenges, a modular architecture consisting of two decoupled main components was designed.
The \acs{medusa} module handles the geometric path planning, while the \acs{kronos} module acts as middleware to form the interface to the industrial robot.

Within the scope of this work, four specific slicing strategies were conceptualized and investigated to varying depths:

\begin{enumerate}
    \item \textbf{Planar Algorithm} Utilizes conventional horizontal slicing planes and serves as a functional reference and validation baseline in the project.
    \item \textbf{Base Algorithm} A fully implemented multi-axis approach.
    It geometrically detects overhangs and dynamically splits the component into a ``trunk'' and ``branches''.
    These branches are printed with their own local growth directions, which enables the support-free manufacturing of overhangs.
    \item \textbf{Sphere Algorithm} A purely conceptual design that decomposes the component into concentric spherical shells around a central starting point.
    It was not implemented, as mechanical weak points are to be expected at the transitions between the spheres due to insufficient bonding areas.
    \item \textbf{Crystal Algorithm} A more complex, field-driven approach.
    It is based on solving the Laplace equation to generate a harmonic scalar field inside the component.
    The print layers follow the natural geometry of the component as curved iso-surfaces, causing the tool orientation to vary continuously.
\end{enumerate}

The evaluation of the overall system was carried out in a simulation using three test geometries (cube, double pyramid, T-piece) in the RViz environment.
The results prove the fundamental feasibility of the established data chain from the \acs{CAD} model to the simulated 6-axis robot movement.
With the Base Algorithm, 90\degree~overhangs could be successfully planned without supports.
The Crystal Algorithm also successfully demonstrates the generation of support-free, fully non-planar toolpaths.

The work identifies the limits of the inverse kinematics solver as well as the future need to integrate real extrusion physics and cross-layer branch collision checking as critical challenges for future development.
Overall, the work provides a valuable and flexibly extensible ``proof of concept'' for closing the software gap in robot-assisted additive manufacturing.
\sig{Ben Stahl}